\documentclass[10pt]{article}

% Packages pour les marges
\usepackage[
    top=2cm,
    bottom=2cm,
    left=2cm,
    right=2cm
]{geometry}

% Police sans serif
% \usepackage{helvet}
% \renewcommand{\familydefault}{\sfdefault}

% Packages existants
\usepackage[french]{babel}
\usepackage[utf8]{inputenc}
\usepackage[T1]{fontenc}
\usepackage{amsmath}
\usepackage{amsfonts}
\usepackage{amssymb}
\usepackage[version=4]{mhchem}
\usepackage{stmaryrd}
\usepackage{listings}

\usepackage{enumitem}
\usepackage{ifthen}
\usepackage{graphicx}
\usepackage{verbatim}
\usepackage{url}
\usepackage{mdframed}
\usepackage[sf]{titlesec}
\titleformat{\section}{\sffamily\large\bfseries}{\thesection}{1em}{}
% Définition de la variable pour afficher les corrections
\newboolean{showSolutions}
% Décommentez la ligne suivante pour afficher les solutions
\input \jobname.adr
% \input \jobname.adr


\title{TD2 - Analyse Numérique Matricielle}

\author{}
\date{}

\newenvironment{solution}
    {\par\vspace{0.5em}\begin{mdframed}[linewidth=0.5pt]\par}
    {\end{mdframed}\par\vspace{0.5em}}

\begin{document}
\sffamily
\maketitle

\section*{A l'examen}
Je vous demanderai :
\begin{itemize}
    \item de décomposer une matrice en un produit $A = LU$ où $L$ est une matrice triangulaire inférieure et $U$ une matrice triangulaire supérieure.
    \item de résoudre un système linéaire en utilisant la décomposition LU.
    \item de calculer l'inverse d'une matrice en utilisant la décomposition LU.
    \item de résoudre un système linéaire en utilisant la méthode de Jacobi.
\end{itemize}

\vspace{1cm}
\section*{Pivot de Gauss et décomposition LU}
On considère les systèmes linéaires suivants :
$$
\left\{\begin{array}{l}
x+y+2 z=3 \\
x+2 y+z=1 \\
2 x+y+z=0
\end{array}\right., \qquad 
\left\{\begin{array}{l}
x+2 z=1 \\
-y+z=2 \\
x-2 y=1
\end{array}\right.
$$

\begin{enumerate}
    \item Ecrivez ces systèmes sous forme matricielle $A X = B$.
    \ifthenelse{\boolean{showSolutions}}{
        \begin{solution}
            Pour le premier système :
            \begin{align*}
                A = \begin{pmatrix}
                    1 & 1 & 2 \\
                    1 & 2 & 1 \\
                    2 & 1 & 1
                \end{pmatrix}
                \quad
                B = \begin{pmatrix}
                    3 \\
                    1 \\
                    0
                \end{pmatrix}
            \end{align*}

            Pour le second système :
            \begin{align*}
                A = \begin{pmatrix}
                    1 & 0 & 2 \\
                    0 & -1 & 1 \\
                    1 & -2 & 0
                \end{pmatrix}
                \quad
                B = \begin{pmatrix}
                    1 \\
                    2 \\
                    1
                \end{pmatrix}
            \end{align*}
        \end{solution}
    }{}


Nous allons suivre un algorithme pour décomposer une matrice $A$ en un produit $A = LU$ où $L$ est une matrice triangulaire inférieure avec des $1$ sur la diagonale et $U$ une matrice triangulaire supérieure.\newline 
\textit{On note souvent $L$ pour Lower, et $U$ pour Upper.}

Cet algorithme est basé sur celui du pivot de Gauss.

    \vspace{1em}
    On rappelle qu'effectuer une action sur les lignes d'une matrice correspond à multiplier à gauche par une matrice $L$: 
    \begin{itemize}
        \item Cette matrice $L$ est la matrice identité à laquelle on a effectué la même opération élémentaire.
        \item L'inverse de cette matrice $L$ est la matrice $L^{-1}$ qui correspond à l'inverse de l'opération élémentaire effectuée.
    \end{itemize}

    \item En suivant l'algorithme du pivot de Gauss, déterminer une matrice $L$ telle que $LA = U$ où $A$ est la matrice du premier système, $L$ une matrice triangulaire inférieure et $U$ une matrice triangulaire supérieure.
    \ifthenelse{\boolean{showSolutions}}{
        \begin{solution}
            La première étape est de faire apparaître des 0 sur la première colonne, en dehors de la diagonale. Les opérations à effectuer sont :
            \begin{itemize}
                \item $L_2 \leftarrow L_2 - L_1$
                \item $L_3 \leftarrow L_3 - 2L_1$
            \end{itemize}
            La matrice $L$ correspondante est :
            \begin{align*}
                L_1 = \begin{pmatrix}
                    1 & 0 & 0 \\
                    -1 & 1 & 0 \\
                    -2 & 0 & 1
                \end{pmatrix}
            \end{align*}
            On obtient la matrice $L_1 A$ :
            \begin{align*}
                L_1 A = \begin{pmatrix}
                    1 & 1 & 2 \\
                    0 & 1 & -1 \\
                    0 & -1 & -3
                \end{pmatrix}
            \end{align*}
            
            On continue avec la deuxième colonne :
            \begin{itemize}
                \item $L_3 \leftarrow L_3 + L_2$
            \end{itemize}
            La matrice $L$ correspondante est :
            \begin{align*}
                L_2 = \begin{pmatrix}
                    1 & 0 & 0 \\
                    0 & 1 & 0 \\
                    0 & 1 & 1
                \end{pmatrix}
            \end{align*}
            On obtient la matrice $L_2 L_1 A$ : 
            \begin{align*}
                L_2 L_1 A = \begin{pmatrix}
                    1 & 1 & 2 \\
                    0 & 1 & -1 \\
                    0 & 0 & -4
                \end{pmatrix}
            \end{align*}
            On a donc $L = L_2 L_1$ et $U = L_2 L_1 A$
        \end{solution}
    }{}

    \item En inversant la matrice $L$ (de manière simple ! vous avez déjà fait tous les calculs nécessaires), déterminer la matrice $L^{-1}$ et écrire $A$ comme un produit $A = L^{-1}U$.\newline
    On appelle cette décomposition, \textbf{la décomposition LU d'une matrice.}
    \ifthenelse{\boolean{showSolutions}}{
        \begin{solution}
            Pour inverser la matrice $L$, on inverse les opérations élémentaires effectuées.

            Inversons tout d'abord $L_1$, les opérations effectuées sont :
            \begin{itemize}
                \item $L_2 \leftarrow L_2 + L_1$
                \item $L_3 \leftarrow L_3 + 2L_1$
            \end{itemize}
            On obtient :    
            \begin{align*}
                L_1^{-1} = \begin{pmatrix}
                    1 & 0 & 0 \\
                    1 & 1 & 0 \\
                    2 & 0 & 1
                \end{pmatrix}
            \end{align*}
            On continue avec $L_2$ :
            \begin{itemize}
                \item $L_3 \leftarrow L_3 - L_2$
            \end{itemize}
            On obtient :
            \begin{align*}
                L_2^{-1} = \begin{pmatrix}
                    1 & 0 & 0 \\
                    0 & 1 & 0 \\
                    0 & -1 & 1
                \end{pmatrix}
            \end{align*}
            On a donc $L^{-1} = L_2^{-1} L_1^{-1}$:
            \begin{align*}
                L^{-1} = \begin{pmatrix}
                    1 & 0 & 0 \\
                    1 & 1 & 0 \\
                    2 & -1 & 1
                \end{pmatrix}
            \end{align*}
            On a donc $A = L^{-1}U$:

            \begin{align*}
                \begin{pmatrix}
                    1 & 1 & 2 \\
                    1 & 2 & 1 \\
                    2 & 1 & 1
                \end{pmatrix}
                = \begin{pmatrix}
                    1 & 0 & 0 \\
                    1 & 1 & 0 \\
                    2 & -1 & 1
                \end{pmatrix}
                \begin{pmatrix}
                    1 & 1 & 2 \\
                    0 & 1 & -1 \\
                    0 & 0 & -4
                \end{pmatrix}
            \end{align*}
        \end{solution}
    }{}

    \item Calculer la décomposition LU de la matrice du deuxième système.
    \ifthenelse{\boolean{showSolutions}}{
        \begin{solution}
            On procède de la même manière que pour le premier système:
            Première opération élémentaire :
            \begin{itemize}
                \item $L_3 \leftarrow L_3 - L_1$
            \end{itemize}
            On obtient :
            \begin{align*}
                L_1 = \begin{pmatrix}
                    1 & 0 & 0 \\
                    0 & 1 & 0 \\
                    -1 & 0 & 1
                \end{pmatrix}
            \end{align*}
            et 
            \[
                L_1 A = \begin{pmatrix}
                    1 & 0 & 2 \\
                    0 & -1 & 1 \\
                    0 & -2 & -2
                \end{pmatrix}
            \]
            On continue avec la deuxième opération élémentaire :
            \begin{itemize}
                \item $L_3 \leftarrow L_3 - 2L_2$
            \end{itemize}
            On obtient :
            \begin{align*}
                L_2 = \begin{pmatrix}
                    1 & 0 & 0 \\
                    0 & 1 & 0 \\
                    0 & -2 & 1
                \end{pmatrix}
            \end{align*}
            et
            \[
                L_2 L_1 A = \begin{pmatrix}
                    1 & 0 & 2 \\
                    0 & -1 & 1 \\
                    0 & 0 & -4
                \end{pmatrix}
            \]
            On a donc $L = L_2 L_1$ et $U = L_2 L_1 A$.

            Pour calculer $L^{-1}$, on inverse les opérations élémentaires effectuées:
            \begin{itemize}
                \item $L_3 \leftarrow L_3 + L_1$
                \item $L_3 \leftarrow L_3 + 2L_2$
            \end{itemize}
            On obtient :
            \begin{align*}
                L^{-1} = \begin{pmatrix}
                    1 & 0 & 0 \\
                    0 & 1 & 0 \\
                    1 & 2 & 1
                \end{pmatrix}
            \end{align*}
            On a donc $A = L^{-1}U$:
            \begin{align*}
                \begin{pmatrix}
                    1 & 0 & 2 \\
                    0 & -1 & 1 \\
                    1 & 2 & 0
                \end{pmatrix}
                = \begin{pmatrix}
                    1 & 0 & 0 \\
                    0 & 1 & 0 \\
                    1 & 2 & 1
                \end{pmatrix}
                \begin{pmatrix}
                    1 & 0 & 2 \\
                    0 & -1 & 1 \\
                    0 & 0 & -4
                \end{pmatrix}
            \end{align*}
        \end{solution}
    }{}



    \vspace{1em}
    La décomposition LU est utilisée dans les logiciels de calcul matriciel pour résoudre les systèmes linéaires, calculer des déterminants et inverser des matrices.
    \item Calculer facilement le déterminant des matrices des deux systèmes ci-dessus.
    \ifthenelse{\boolean{showSolutions}}{
        \begin{solution}
            Le déterminant d'un produit de matrices est le produit des déterminants et le déterminant d'une matrice triangulaire est le produit des éléments diagonaux.
            Donc pour le premier système, on a :
            \[
                \det(A) = \det(L) \det(U) = 1 \times (-4) = -4
            \]
            et pour le second système, on a :
            \[
                \det(A) = \det(L) \det(U) =  1 \times 4 = 4
            \]
        \end{solution}
    }{}

    \item Résoudre le premier système linéaire en utilisant la décomposition LU, on pourra poser la variable intermédiaire $Y = U X$.

    \ifthenelse{\boolean{showSolutions}}{
        \begin{solution}
            On résout le premier système notons 
            \[
            L  = \begin{pmatrix}
                1 & 0 & 0 \\
                1 & 1 & 0 \\
                2 & -1 & 1
            \end{pmatrix}
            \]
            et
            \[
                U = \begin{pmatrix}
                    1 & 1 & 2 \\
                    0 & 1 & -1 \\
                    0 & 0 & -4
                \end{pmatrix}
            \]
            On veut résoudre $LU X = B$ avec $B = \begin{pmatrix} 3 \\ 1 \\ 0 \end{pmatrix}$.

            On pose $Y = U X$ et on résout $LY = B$ :
            \[
                \begin{pmatrix}
                    1 & 0 & 0 \\
                    1 & 1 & 0 \\
                    2 & -1 & 1
                \end{pmatrix}
                \begin{pmatrix}
                    y_1 \\
                    y_2 \\
                    y_3
                \end{pmatrix}
                = \begin{pmatrix} 3 \\ 1 \\ 0 \end{pmatrix}
                \]
                C'est un système triangulaire : 
                \begin{align*}
                    y_1 &= 3 \\
                    y_2 &= 1 - y_1 = 1 - 3 = -2 \\
                    y_3 &= 0 + y_2 - 2y_1 = 0 + (-2) - 6 = -8
                \end{align*}
                On a donc $Y = \begin{pmatrix} 3 \\ -2 \\ -8 \end{pmatrix}$.

                On résout maintenant $U X = Y$ :
                \[
                    \begin{pmatrix}
                        1 & 1 & 2 \\
                        0 & 1 & -1 \\
                        0 & 0 & -4
                    \end{pmatrix}
                    \begin{pmatrix}
                        x_1 \\
                        x_2 \\
                        x_3
                    \end{pmatrix}
                    = \begin{pmatrix} 3 \\ -2 \\ -8 \end{pmatrix}
                \]
                C'est un système triangulaire : 
                \begin{align*}
                    x_3 &= \frac{-8}{-4} = 2 \\
                    x_2 &= -2 + x_3 = -2 + 2 = 0 \\
                    x_1 &= 3 - x_2 - 2x_3 = 3 - 0 - 4 = -1
                \end{align*}
                On a donc $X = \begin{pmatrix} -1 \\ 0 \\ 2 \end{pmatrix}$.

                On peut vérifier que ce vecteur est bien solution du système initial.
        \end{solution}
    }{}

    \item Soit $A$ la matrice du premier système linéaire, calculer l'inverse de $A$ en utilisant la décomposition LU.\newline
    Pour cela, on cherche une matrice inconnue $M$ telle que $A M = I$ où $I$ est la matrice identité.
    \begin{itemize}
        \item On pose $N = UM$, expliquer comment déterminer les coefficients de la matrice $N$. 
        \item Expliquer comment déterminer les coefficients de la matrice $M$.
    \end{itemize}
    \ifthenelse{\boolean{showSolutions}}{
        \begin{solution}
            On cherche une matrice $M$ telle que $A M = I$.
            On pose $N = UM$ et on résout $L N = I$, on note $N = (n_{ij})$ et $M = (m_{ij})$.
            
            Il faut résoudre $L N = I$ : 
            \[
                \begin{pmatrix}
                    1 & 0 & 0 \\
                    1 & 1 & 0 \\
                    2 & -1 & 1
                \end{pmatrix}
                \begin{pmatrix}
                    n_{11} & n_{12} & n_{13} \\
                    n_{21} & n_{22} & n_{23} \\
                    n_{31} & n_{32} & n_{33}
                \end{pmatrix}
                = \begin{pmatrix} 1 & 0 & 0 \\ 0 & 1 & 0 \\ 0 & 0 & 1 \end{pmatrix}
            \]
            Cette équation matricielle fournit 3 systèmes linéaires triangulaires, le premier est donné par la première colonne des matrices : 
            \begin{align*}
                n_{11} &= 1 \\
                n_{11} + n_{21} &= 0 \\
                2 n_{11} - n_{21} + n_{31} &= 0
            \end{align*}
            On obtient donc $n_{11} = 1$, $n_{21} = -1$ et $n_{31} = -3$.
            
            On continue avec la deuxième colonne :
            \begin{align*}
                n_{12} &= 0 \\
                n_{12} + n_{22} &= 1 \\
                2n_{12} - n_{22} + n_{32} &= 0
            \end{align*}
            On obtient donc $n_{12} = 0$, $n_{22} = 1$ et $n_{32} = 1$.
            
            On continue avec la troisième colonne :
            \begin{align*}
                n_{13} &= 0 \\
                n_{13} + n_{23} &= 0 \\
                2n_{13} - n_{23} + n_{33} &= 1
            \end{align*}
            On obtient donc $n_{13} = 0$, $n_{23} = 0$ et $n_{33} = 1$.

            On a donc $N = \begin{pmatrix} 1 & 0 & 0 \\ -1 & 1 & 0 \\ -3 & 1 & 1 \end{pmatrix}$. C'est bien l'inverse de la matrice $L$. 

            On cherche maintenant $M$ tel que $U M = N$ : 
            \[
                \begin{pmatrix}
                    1 & 1 & 2 \\
                    0 & 1 & -1 \\
                    0 & 0 & -4
                \end{pmatrix}
                \begin{pmatrix}
                    m_{11} & m_{12} & m_{13} \\
                    m_{21} & m_{22} & m_{23} \\
                    m_{31} & m_{32} & m_{33}
                \end{pmatrix}
                = \begin{pmatrix} 1 & 0 & 0 \\ -1 & 1 & 0 \\ -3 & 1 & 1 \end{pmatrix}
            \]
            Cette équation matricielle fournit à nouveau 3 systèmes linéaires triangulaires, le premier est donné par la première colonne des matrices : 
            \begin{align*}
                m_{11} + m_{21} + 2m_{31} &= 1 \\
                m_{21} - m_{31} &= -1 \\
                -4 m_{31} &= -3
            \end{align*}
            On obtient donc $m_{11} = \frac{-1}{4}$, $m_{21} = \frac{-1}{4}$ et $m_{31} = \frac{3}{4}$.

            On continue avec la deuxième colonne :
            \begin{align*}
                m_{12} + m_{22} + 2m_{32} &= 0 \\
                m_{22} - m_{32} &= 1 \\
                -4 m_{32} &= 1
            \end{align*}
            On obtient donc $m_{12} = \frac{-1}{4}$, $m_{22} = \frac{3}{4}$ et $m_{32} = \frac{-1}{4}$.

            On finit avec la troisième colonne :
            \begin{align*}
                m_{13} + m_{23} + 2m_{33} &= 0 \\
                m_{23} - m_{33} &= 0 \\
                -4 m_{33} &= 1
            \end{align*}
            On obtient donc $m_{13} = \frac{3}{4}$, $m_{23} = \frac{-1}{4}$ et $m_{33} = \frac{-1}{4}$.

            Finalement,  
            \[
                M = \frac{1}{4} \begin{pmatrix} -1 & -1 & 3 \\ -1 & 3 & -1 \\ 3 & -1 & -1 \end{pmatrix}
            \]
            On peut vérifier que $A M = I$.
        \end{solution}
    }{}
    \item A votre avis, pourquoi cette méthode est-elle privilégiée par rapport au calcul de $L^{-1}, U^{-1}$ puis du produit $U^{-1}L^{-1}$ ?
    \ifthenelse{\boolean{showSolutions}}{
        \begin{solution}
            Cette méthode économise l'espace en mémoire, il faut imaginer des grandes matrices 1000x1000, le calcul de l'inverse de $A$ est très coûteux en temps et en mémoire.
        \end{solution}
    }{}

    \item Quelle est la complexité de l'algorithme de décomposition LU ?
    \ifthenelse{\boolean{showSolutions}}{
        \begin{solution}
            La complexité de l'algorithme de décomposition LU est le même que le pivot de Gauss : $O(n^3)$.
        \end{solution}
    }{}
\end{enumerate}


\section*{Méthode de Jacobi}
Si les systèmes linéaires sont trop grands, on utilise une méthode différente pour la résolution : la méthode de Jacobi.

Une matrice $A$ d'un système linéaire $A X=B$ peut toujours s'écrire sous la forme $A=D-E-F$ avec
\begin{itemize}
    \item[$\bullet$] $D$ la matrice composée de la diagonale de $A$.
    \item[$\bullet$] $-E$ matrice triangulaire inférieure composée des éléments de $A$ en-dessous de la diagonale.
    \item[$\bullet$] $-F$ matrice triangulaire supérieure composée des éléments de $A$ au-dessus de la diagonale.
\end{itemize}

\begin{enumerate}
    \item Montrer que $X$ vérifie 
    \[
        X = D^{-1}B + D^{-1}(E + F) X\,.
    \]
    \ifthenelse{\boolean{showSolutions}}{
        \begin{solution}
            On a $A X = B$ donc $D X = B + E X + F X$.
            On en déduit que $X = D^{-1}B + D^{-1}(E + F) X$.
        \end{solution}
    }{}

    \vspace{1em}
    L'idée de la méthode de Jacobi est de construire une suite $\Big(X^{(k)}\Big)_{k \in \mathbb{N}}$ qui devrait converger vers la solution $X$ du système linéaire.
    On part d'une valeur initiale $X^{(0)} = D^{-1}B$ et on construit la suite par récurrence :
    \[
        X^{(k+1)} = D^{-1}B + D^{-1}(E + F) X^{(k)}\,.
    \]
    \item Lorsqu'on travaille avec une suite réelle $(u_n)$ arithmético-géométrique : 
    \[
        u_{n+1} = a + b u_n\,,
    \]
    rappelez pourquoi la suite converge vers une solution $u^*$ si et seulement si $|b| < 1$.\newline
    Pour cela, on suppose que la suite converge, on note $u^*$ sa limite et on montre que la suite $(u_{n+1} - u^*)$ est géométrique.

    \ifthenelse{\boolean{showSolutions}}{
        \begin{solution}
            En notant $u^*$ la limite de la suite (en supposant qu'elle existe), on a $u^* = a + b u^*$.
            On en déduit que $u^* = \frac{a}{1-b}$.
            La suite récurrente se réécrit alors 
            \[
                u_{n+1} - u^* = b (u_n - u^*).
            \]
            La suite géométrique $(u_n - u^*)$ converge si et seulement si $|b| < 1$.
        \end{solution}
    }{}

    \item Dans la méthode de Jacobi, la suite récurrente est une suite de matrice. A votre avis, quel est le critère sur la matrice $D^{-1}(E + F)$ pour que la méthode de Jacobi converge ?
    \textit{Pensez à la diagonalisation.}
    \ifthenelse{\boolean{showSolutions}}{
        \begin{solution}
            La méthode de Jacobi converge si et seulement si la plus grande valeur propre de la matrice $D^{-1}(E + F)$ est inférieure strictement à 1.
        \end{solution}
    }{}

    \vspace{1em}
     On veut résoudre par la méthode itérative de Jacobi le système suivant :
$$
\left\{\begin{array}{l}
4 x+2 y=64 \\
x+4 y+2 z=84 \\
y+5 z=77
\end{array}\right.
$$

    \item Définir les matrices $D$ et $B$ qui interviennent dans la méthode de Jacobi.
    \ifthenelse{\boolean{showSolutions}}{
        \begin{solution}
            On a $D = \begin{pmatrix} 4 & 0 & 0 \\ 0 & 4 & 0 \\ 0 & 0 & 5 \end{pmatrix}$ et $B = \begin{pmatrix} 64 \\ 84 \\ 77 \end{pmatrix}$.
        \end{solution}
    }{}

    \item Calculer $X^{(0)}=D^{-1} B$.
    \ifthenelse{\boolean{showSolutions}}{
        \begin{solution}
            On a $X^{(0)} = \begin{pmatrix} 16 \\ 21 \\ 77/5 \end{pmatrix}$.
        \end{solution}
    }{}

    \item Calculer $X^{(1)}$.
    \ifthenelse{\boolean{showSolutions}}{
        \begin{solution}
            Calculons $D^{-1}(E + F)$ :
            \begin{align*}
                D^{-1}(E + F) &= \begin{pmatrix} \frac{1}{4} & 0 & 0 \\ 0 & \frac{1}{4} & 0 \\ 0 & 0 & \frac{1}{5} \end{pmatrix} \begin{pmatrix} 0 & -2 & 0 \\ -1 & 0 & -2 \\ 0 & -1 & 0 \end{pmatrix} \\
                &= \frac{-1}{20} \begin{pmatrix} 0 & 10 & 0 \\ 5 & 0 & 10 \\ 0 & 4 & 0 \end{pmatrix}
            \end{align*}
            On en déduit que 
            \[
                X^{(1)} = \begin{pmatrix} 16 \\ 21 \\ 77/5 \end{pmatrix} + \frac{-1}{20} \begin{pmatrix} 0 & 10 & 0 \\ 5 & 0 & 10 \\ 0 & 4 & 0 \end{pmatrix} \begin{pmatrix} 16 \\ 21 \\ 77/5 \end{pmatrix}
            \]
            % On obtient $X^{(1)} = \begin{pmatrix} 16 \\ 21 \\ 77/5 \end{pmatrix}$.

            % On a \[
            % X^{(1)} = \begin{pmatrix} 16 \\ 21 \\ 77/5 \end{pmatrix} + \begin{pmatrix} 0 & -2 & 0 \\ -1 & 0 & -2 \\ 0 & -1 & 0 \end{pmatrix} \begin{pmatrix} 16 \\ 21 \\ 77/5 \end{pmatrix}
            % \]
            
        \end{solution}
    }{}

    \item Calculer la solution exacte en factorisant $A$ en $LU$ et en résolvant les deux systèmes triangulaires, puis l'erreur commise par la solution approchée $X^{(1)}$.
    % \ifthenelse{\boolean{showSolutions}}{
    %     \begin{solution}
    %         La solution exacte est $X = \begin{pmatrix} 16 \\ 21 \\ 77/5 \end{pmatrix}$.
    %         L'erreur commise est $X - X^{(1)} = \begin{pmatrix} 0 \\ 0 \\ 0 \end{pmatrix}$.
    %     \end{solution}
    % }{}
    \vspace{1em}

    Cette méthode converge toujours si la matrice $A$ vérifie l'un des critères suivants : 
    \begin{itemize}
        \item[$\bullet$] est à diagonale strictement dominante, c'est-à-dire que pour tout $i$, $\sum_{j \neq i} |a_{ij}| < |a_{ii}|$.
        \item[$\bullet$] est symétrique définie positive.
    \end{itemize}
    \item Montrer le premier critère. Généralisez à des matrices $n \times n$.
\end{enumerate}



\end{document}
